\section{The rank-one obstruction}\label{sec:rank}

We prove an upper bound for every knot obtained by following \(U\)
with a braid in the subgroup acting on \eqref{eq:framed}. That
subgroup fixes the first and fourth meridian words. These two fixed
positions suffice for the rank-one argument.

\begin{lemma}\label{lem:spans}
Suppose that \(\tau\) belongs to
\(\langle\sigma_2,g,\sigma_5,\ldots,\sigma_{b-1}\rangle\) and
\(\cl{U\tau}\) is a knot. Let its group act on a finite-dimensional
space \(V\) over a field so that each meridian has displacement
of rank one. Write each meridian as \(M=I-u f\), where \(u\in V\)
and \(f\in V^*\). Then at least one of the spans of all the \(u\)'s
and all the \(f\)'s has dimension at most \(b-1\), and both
dimensions are at most \(b\).
\end{lemma}

\begin{proof}
Each of the two spans is already spanned by its \(b\) top-strand
vectors or covectors. The root span of the generating meridians is
invariant under those generators and their inverses, and the
covector span is invariant under right multiplication by the same
matrices. Conjugating a generator therefore introduces neither a
new root direction nor a new covector direction.

Suppose both spans have dimension \(b\). In particular both
families of four vectors and covectors on the first four strands
are separately independent. Set
\[
 \mu=\det M\ne0,\qquad f(u)=1-\mu,\qquad
 M^{-1}=I+\mu^{-1}u f.
\]
The first two identities hold for every meridian, with one common
\(\mu\). When \(M_a=I-u_a f_a\), conjugation by a matrix \(J\)
replaces its pair by \((Ju_a,f_aJ^{-1})\). Two pairs for the same
nonzero displacement differ by reciprocal scaling.

Here is the crossing trace through \(U\). We use arc subscripts on
the matrices \(M_i\) in this proof. A row \((o,a,c,\epsilon)\)
represents \(M_c=M_o^\epsilon M_aM_o^{-\epsilon}\).
The input labels are \((1,0,8,13)\) and the output labels are
\((1,6,16,13)\).
\begin{equation}\label{eq:prefix}
\begin{array}{c|rrrr@{\qquad}c|rrrr}
0&0&1&2&-&1&8&2&3&-\\
2&3&13&14&+&3&14&3&4&+\\
4&4&8&9&-&5&9&14&15&+\\
6&15&9&10&+&7&10&4&5&-\\
8&10&0&1&-&9&1&10&11&-\\
10&5&11&12&-&11&12&5&6&-\\
12&6&12&13&-&13&13&15&16&+
\end{array}
\end{equation}
The repeated outer labels express closure, since \(\tau\) fixes
the first and fourth entries. Each row is obtained by \(T\) or
\(T^{-1}\) in the order of the letters in \(U\).

Put \(A=M_2,\ B=M_8,\ C=M_{13},\ D=M_0\).
Replacing \(M_1\) by \(A=D^{-1}M_1D\) changes the root by a
multiple of the root of \(D\) and the covector by a multiple of
its covector, preserving both independent families.
Let \(H\) and \(F\) denote the root and covector spans of
\(A,B,C\), respectively. Both have dimension three, and in the
pair \((v,\nu)\) of \(D\) one has \(v\notin H\) and \(\nu\notin F\).
Rows \(1\) through \(7\) and row \(13\) propagate pairs in \(H,F\).

Rows \(0\) and \(8\) give \(DAD^{-1}=X^{-1}DX\) for
\(X=M_{10}\). If \(A,X\) have pairs \((a,f)\) and \((x,\xi)\)
in \(H,F\), comparing the pairs on both sides gives
\[
 \begin{aligned}
 a-\nu(a)v &= \lambda\bigl(v+\mu^{-1}\xi(v)x\bigr),\\
 f+\mu^{-1}f(v)\nu &= \lambda^{-1}\bigl(\nu-\nu(x)\xi\bigr)
 \end{aligned}
\]
for some \(\lambda\ne 0\). The coefficients on \(v,\nu\)
give \(\lambda=-\nu(a)\) and \(f(v)=\mu\lambda^{-1}\).
The other components give
\(a=\lambda\mu^{-1}\xi(v)x\) and \(f=-\lambda^{-1}\nu(x)\xi\).
Applying \(\nu\) to the expression for \(a\) gives
\(\nu(x)\xi(v)=-\mu\). Thus \(af=x\xi\) and \(X=A\).
The relation is now \(DAD=ADA\).
Row \(9\), together with \(M_1=DAD^{-1}\), gives \(M_{11}=D\).

Next, rows \(10\)--\(12\) send \((D,P)\) to \((P,N)\),
then \((N,Q)\), then \((Q,C)\), where \(P=M_5,\ N=M_{12}\)
and \(Q=M_6\). Write the pair of \(P\) as \((p,\pi)\) in
\(H,F\), and put \(x_1=\pi(v),\ y_1=\nu(p),\
h_1=1+x_1y_1/\mu\). Conjugating at each step gives
\[
\begin{array}{c|cc}
 &\text{root}&\text{covector}\\ \hline
N&v+(x_1/\mu)p&\nu-y_1\pi\\
Q&y_1v+h_1p&-(x_1/\mu)\nu+h_1\pi\\
C&h_1v+(x_1/\mu)(1+h_1)p&h_1\nu-y_1(1+h_1)\pi.
\end{array}
\]
For example \(N=P^{-1}DP\); the next row uses \(Q=N^{-1}PN\).
The last pair lies in \(H,F\), so \(h_1=0\) and
\(x_1y_1=-\mu\). Its remaining coefficients multiply to one,
which gives \(P=C\).

Of the eight rows supported in \(H,F\), rows \(1\)--\(7\)
suffice for the contradiction. They now satisfy the hypotheses
of Lemma~\ref{lem:seven} and contain the three independent pairs
of \(A,B,C\). This proves the span bound.
\end{proof}

\begin{lemma}\label{lem:seven}
Suppose the rank-one meridian matrices of common nonzero
determinant \(\mu\) satisfy the seven signed relations in
rows \(1\)--\(7\) of \eqref{eq:prefix}, with
\(M_{10}=M_2\) and \(M_5=M_{13}\). Their root and covector
spans cannot both have dimension three.
\end{lemma}

\begin{proof}
Reorder the crossings so that the propagation comes first:
\[
\begin{gathered}
 (9,14,15,+),\ (15,9,2,+),\ (8,2,3,-),\
 (14,3,4,+),\ (3,14,13,-),\\
 (4,9,8,+),\qquad (2,13,4,+).
\end{gathered}
\]
Here the rows with over-arcs \(3,4,2\) and output arcs
\(13,8,4\), respectively, solve the corresponding prefix relations
for the other under-arc. The first five crossings propagate all the
matrices from \(M_8,M_{14},M_9\). If both spans have dimension three,
the corresponding pairs \((u,f),(v,g_0),(w,j)\)
are separately independent. Set
\[
\begin{gathered}
 x=f(v),\quad X=g_0(u),\quad y=f(w),\quad Y=j(u),\quad
 z=g_0(w),\quad Z=j(v),\qquad \eta=1-\mu,\\
 h=1+zZ/\mu,\qquad t=-zx/\mu+hy,\qquad \ell=ZX+hY,\\
 \kappa=tX/\mu+z(h-\eta/\mu),\qquad
 \Lambda=-\ell x+(\eta+h)Z .
\end{gathered}
\]
In the bases \((u,v,w)\) and \((f,g_0,j)\), propagation gives
\[
\begin{array}{c|cc}
 i&u_i&f_i\\\hline
 2&(0,-z/\mu,h)&(0,Z,h)\\
 3&(t/\mu,-z/\mu,h)&(-\ell,Z,h)\\
 4 &(t/\mu,-z/\mu-\kappa,h)&(-\ell,Z+\Lambda/\mu,h)\\
 13&(\Lambda t/\mu^2,1-\Lambda z/\mu^2,\Lambda h/\mu)
       &(\kappa\ell,1-\kappa Z,-\kappa h).
\end{array}
\]
For the first row, \(M_{15}=M_{9}M_{14}M_{9}^{-1}\) has pair
\((v-Zw,g_0+(z/\mu)j)\), and \(M_{2}=M_{15}M_{9}M_{15}^{-1}\).
The next two rows conjugate by \(M_{8}^{-1}\) and \(M_{14}\),
respectively, with \(t=f(u_{2}),\ \ell=f_{2}(u)\) and
\(\kappa=g_0(u_{3}),\ \Lambda=f_{3}(v)\).
The last row uses \(M_{13}=M_{3}^{-1}M_{14}M_{3}\).

At the crossing \((4,9,8,+)\), put \(P=f_{4}(w)\) and
\(Q=j(u_{4})\). The first and third coefficients of
\((M_{4}w,jM_{4}^{-1})=(\gamma u,\gamma^{-1}f)\) give
\[
 -Pt/\mu=\gamma,\qquad Ph=1,\qquad
 -Q\ell/\mu=\gamma^{-1},\qquad Qh/\mu=-1.
\]
Thus \(P,Q,t,\ell,h\) are nonzero. The middle coefficients
give \(\kappa=-z/\mu\) and \(\Lambda=-\mu Z\).
At \((2,13,4,+)\) the \(M_{2}\) pair has zero first
coefficients. Comparing the first coefficients of the pairs
of \(M_{13}\) and \(M_{2}^{-1}M_{4}M_{2}\) therefore gives
\[
 \Lambda t/\mu^2=\beta t/\mu,\qquad
 \kappa\ell=-\beta^{-1}\ell
\]
with \(\beta\ne0\). Cancelling \(t,\ell\) yields
\(zZ=\kappa\Lambda=-\mu\). This makes \(h=0\), contrary to
\(Ph=1\).
\end{proof}

For an irreducible representation in Cornwell's correspondence,
the root span is nonzero and invariant, hence equals \(V\).
The common kernel of the covectors is the fixed subspace and
is zero, so the covectors span \(V^*\).
Lemma~\ref{lem:spans} bounds the dimension by \(b-1\) and
proves the augmentation inequality in Theorem~\ref{thm:main}.
